\documentclass[12pt,a4paper]{article}
\input{../pri-end/T0_preamble_standalone_De}

\hypersetup{
	pdftitle={181 · Begründung der 4D-Torusgeometrie und Herleitung von L0},
	pdfauthor={Johann Pascher},
	pdfsubject={FFGFT: Warum der 4D-Torus aus dem Lagrangian folgt und L0 die kleinste stabile Windung ist}
}

\fancyhead[L]{\small\textcolor{t0gray}{FFGFT · 4D-Torusgeometrie}}
\fancyhead[R]{\small\textcolor{t0gray}{Johann Pascher · 2026}}

\title{%
	{\LARGE\bfseries\color{t0blue}
	 Begründung der 4D-Torusgeometrie}\\[0.5em]
	{\large und Herleitung von $L_0$ als kleinste stabile Windung}\\[1em]
	{\normalsize Johann Pascher \quad $\cdot$ \quad März 2026}\\[0.5em]
	{\small \url{https://github.com/jpascher/T0-Time-Mass-Duality}}
}
\date{März 2026}

\begin{document}
	\maketitle
	\tableofcontents
	\newpage

	\section*{Abstract}

Der 4D-Torus ist kein Postulat der T0-Theorie. Er folgt zwingend aus drei unabhängigen Bedingungen: (1) Divergenzfreiheit der Quantenfeldtheorie verlangt periodische Topologie, (2) der Wert $\xipar = 4/30000$ kodiert genau vier Raumzeit-Dimensionen, (3) stabile Windungszahlen auf dem Torus erzeugen die beobachteten Quantenzahlen. Die minimale Längenskala $L_0 = \xipar \cdot \lP$ ist geometrisch die Länge einer einzelnen Torusgitterzelle — der kleinsten stabilen Windung. Lagrangian und Torus liefern $L_0$ konsistent auf zwei unabhängigen Wegen.

\section{Warum eine vierte Dimension?}

\subsection{Das Divergenzproblem der Quantenfeldtheorie}

Die Quantenfeldtheorie leidet an einem fundamentalen Problem: Die Vakuumenergie divergiert:
\begin{equation}
	E_{\text{vak}} = \sum_n \frac{1}{2}\hbar\omega_n \rightarrow \infty
\end{equation}
Die beobachtete kosmologische Konstante ist $10^{120}$-mal kleiner als die QFT-Vorhersage. Das Standardmodell löst dieses Problem durch künstliche Renormierung. T0/FFGFT löst es geometrisch.

\subsection{Periodische Geometrie als Lösung}

Auf einer unendlichen Linie divergiert die harmonische Reihe. Auf einem Kreis konvergiert sie — Euler 1735:
\begin{equation}
	\sum_{n=1}^{\infty} \frac{1}{n^2} = \frac{\pi^2}{6}
\end{equation}

\begin{keypoint}[Topologischer Schluss]
Die periodische Geometrie unterdrückt hohe Moden automatisch. Wenn die Raumzeit eine periodische Topologie hat, verschwinden UV-Divergenzen geometrisch — ohne künstliche Renormierung. Das $\pi^2$ in Eulers Ergebnis ist nicht Zahlentheorie, sondern die Signatur der Topologie.
\end{keypoint}

\subsection{Vier Dimensionen aus \texorpdfstring{$\xi$}{xi}}

Der Parameter $\xipar = 4/30000$ kodiert die Dimensionalität der Raumzeit:
\begin{equation}
	\xipar_d \sim (10^{-1})^d \quad \Rightarrow \quad \text{für } d=4: \quad \xipar_4 \sim 10^{-4}
\end{equation}

\begin{keyresult}[Vier Dimensionen aus \texorpdfstring{$\xi$}{xi}]
Die \textbf{4} im Zähler von $\xipar = 4/30000$ steht direkt für die vier Raumzeit-Dimensionen. Die vierte Dimension ist nicht postuliert — sie ist im Wert von $\xipar$ bereits enthalten.
\end{keyresult}

\section{Warum ein Torus?}

\subsection{Windungszahlen als Erhaltungsgrößen}

Auf einem Torus sind Windungszahlen topologisch erhalten. Sie erzeugen stabile diskrete Zustände:
\begin{align}
	\text{Spin} &= \text{Windungszahl } w = n_\phi / n_\theta \\
	\text{Ladung} &= \text{Fluss } \Phi = n \cdot h/e \\
	\text{Massen} &= \text{Torus-Obertonfrequenzen}
\end{align}

\begin{keypoint}[Torus vs.\ Kugel]
Eine Kugel $S^3$ ist einfach zusammenhängend — keine stabilen Windungszahlen. Der Torus ist nicht einfach zusammenhängend — er erzeugt genau die beobachteten Quantenzahlen. Das ist der topologische Grund für die Toruswahl.
\end{keypoint}

\subsection{Torus-Eigenwerte und Teilchenmassen}

Der Laplace-Operator auf einem $d$-dimensionalen Torus mit Radien $R_i$ hat Eigenwerte:
\begin{equation}
	E_{\vec{n}} = \frac{\hbar^2}{2} \sum_i \frac{n_i^2}{R_i^2}, \qquad \vec{n} \in \mathbb{Z}^d
\end{equation}
Die Teilchenmassen erscheinen als Torusobertonfrequenzen mit Schrittweite $1/3 = 1/D_f$:
\begin{equation}
	m_n = m_0 \cdot n^{1/3}
\end{equation}
Daraus folgt die Koide-Formel für die Leptonmassen als geometrische Konsequenz.

\subsection{Kompaktifizierung}

Die vierte Dimension ist auf den Radius $r_4$ aufgerollt. Dieser Radius muss mit der Sub-Planck-Skala aus dem Lagrangian übereinstimmen — sonst wäre die Theorie inkonsistent. Die Bedingung:
\begin{equation}
	r_4 = L_0 = \xipar \cdot \lP
\end{equation}
Die Topologie ist daher:
\begin{equation}
	\mathbb{R}^3 \times S^1 \quad \text{mit} \quad r_4 = \xipar \cdot \lP
\end{equation}

\section{Herleitung von \texorpdfstring{$L_0$}{L0} aus der Torusgeometrie}

\subsection{Die Windungsdichte}

Der Torus hat $f = 1/\xipar$ Sub-Planck-Zellen pro Planck-Länge. Mit $\xipar = 4/30000$:
\begin{equation}
	f = \frac{1}{\xipar} = \frac{30000}{4} = 7500 \quad \text{Zellen pro } \lP
\end{equation}
Diese Zahl folgt direkt aus der geometrischen Windungsdichte $30000/4$ des Torus. Sie ist nicht willkürlich — die Primfaktorzerlegung $7500 = 2^2 \times 3 \times 5^4$ kodiert die Symmetriestruktur der Raumzeit.

\subsection{Kleinste stabile Windung}

Die kleinste stabile Windung auf dem Torus hat die Länge einer einzelnen Gitterzelle:
\begin{equation}
	L_0 = \frac{\lP}{f} = \lP \cdot \xipar
\end{equation}
Numerisch:
\begin{equation}
	L_0 = 1{,}616 \times 10^{-35}\,\text{m} \times \frac{4}{30000} = 1{,}616 \times 10^{-35}\,\text{m} \times 1{,}333 \times 10^{-4}
\end{equation}

\begin{keyresult}[Minimale Längenskala aus Torusgeometrie]
\begin{equation}
	\boxed{L_0 = \xipar \cdot \lP = 2{,}155 \times 10^{-39}\,\text{m}}
\end{equation}
$L_0$ ist die Länge einer einzelnen Torusgitterzelle — der kleinsten stabilen Windung.
\end{keyresult}

\subsection{Geometrische Interpretation}

$L_0$ ist nicht einfach $\xipar$ mal die Planck-Länge — $L_0$ ist die Länge einer einzelnen Torusgitterzelle. Unterhalb von $L_0$:
\begin{itemize}
	\item Verliert das Konzept „Abstand" seine physikalische Bedeutung.
	\item Sind alle Vakuumfluktuationen vollständig wirksam.
	\item Wird die fraktale Korrektur $\Kfrak$ strukturell notwendig.
\end{itemize}

\subsection{Konsistenz der zwei Ableitungswege}

$L_0$ wird nicht doppelt hergeleitet — Lagrangian und Torus liefern dieselbe Zahl auf zwei unabhängigen Wegen:

\begin{keyresult}[Zwei konsistente Wege zu \texorpdfstring{$L_0$}{L0}]
\begin{align}
	\text{Lagrangian:} \quad & g_T = \xipar \cdot m \;\Rightarrow\; \xipar = 4/30000 \notag \\
	& \Downarrow \notag \\
	\text{Torus:} \quad & f = 1/\xipar = 7500 \text{ Zellen pro } \lP \notag \\
	& \Downarrow \notag \\
	\text{Kleinste Windung:} \quad & L_0 = \lP / f = \xipar \cdot \lP \notag \\
	& \Downarrow \notag \\
	& \boxed{L_0 = 2{,}155 \times 10^{-39}\,\text{m}} \quad \text{(konsistent aus beiden Wegen)}
\end{align}
Der Lagrangian liefert $\xipar$. Der Torus erklärt was $\xipar$ geometrisch bedeutet: die reziproke Windungsdichte.
\end{keyresult}

\section{Abgrenzung zu anderen Ansätzen}

\subsection{Stringtheorie}

Stringtheorie benötigt 10 oder 11 Dimensionen. Der T0-Lagrangian mit einem einzigen Parameter $\xipar$ reicht mit 4 aus. Jede zusätzliche Dimension wäre ein Postulat — in FFGFT gibt es keine freien Parameter.

\subsection{Schleifen-Quantengravitation}

LQG diskretisiert den Raum in Spin-Netzwerke. T0 diskretisiert die Raumzeit durch die periodische Torusgeometrie. Der Unterschied: In T0 ist die Diskretheit eine topologische Konsequenz, nicht ein zusätzliches Postulat.

\subsection{Asymptotische Sicherheit}

AS sucht einen UV-Fixpunkt als Ergebnis des Renormierungsgruppen-Flusses. In T0 gibt es keinen Renormierungsgruppen-Fluss — $\xipar$ ist eine geometrische Konstante, kein laufender Parameter. Der Fixpunkt ist nicht gesucht sondern von Anfang an geometrisch festgelegt.

\section{Physikalische Konsequenzen}

\subsection{Vakuumkatastrophe gelöst}

Auf dem fraktalen Torus konvergiert die Vakuumenergie durch periodische Unterdrückung hoher Frequenzen. Keine künstliche Renormierung nötig.

\subsection{Drei Generationen als Obertonsstruktur}

Elektron, Myon und Tau sind nicht drei unabhängige Teilchen — sie sind Grundton und erste Obertöne eines einzigen Torus-Resonators. Die Hierarchie folgt aus den Torus-Eigenwerten:
\begin{equation}
	m_e : m_\mu : m_\tau \approx 1 : 206{,}8 : 3477
\end{equation}

\subsection{Kein Urknall und keine Frequenzabhängigkeit}

Der Torus hat keine Ränder. Die beobachtete Rotverschiebung entsteht nicht durch Expansion, sondern durch geometrischen Energieverlust — die fraktale Summation verlängert die effektive Wegstrecke von Photonen durch das T0-Vakuumfeld:
\begin{equation}
	z \approx \xipar \cdot \ln\!\left(\frac{d}{\lP}\right)
\end{equation}

\begin{keypoint}[Frequenzunabhängigkeit der Rotverschiebung]
Da die Wegverlängerung eine Eigenschaft der Raumzeit-Geometrie selbst ist, nehmen rote und blaue Photonen, die am selben Ort starten, exakt dieselbe verlängerte Wegstrecke. Ihre Wellenlängen werden daher prozentual gleich gestreckt — es gibt keine Frequenzabhängigkeit. Dies unterscheidet den T0-Mechanismus fundamental von einfachen ``Tired Light''-Modellen, die an der beobachteten Frequenzunabhängigkeit scheitern. Ausführlich hergeleitet in Dokument 026 (Geometrische Kosmologie).
\end{keypoint}

\section{Zusammenfassung}

\begin{keyresult}[Drei unabhängige Bedingungen zeigen auf denselben Torus]
\begin{align}
	(1)\; & \text{Divergenzfreiheit} \;\Rightarrow\; \text{periodische Topologie} \notag \\
	(2)\; & \xipar = 4/30000 \;\Rightarrow\; \text{genau 4 Dimensionen} \notag \\
	(3)\; & \text{Windungszahlen} \;\Rightarrow\; \text{stabile Quantenzahlen} \notag
\end{align}
Alle drei zeigen auf denselben Schluss:
\begin{equation}
	\boxed{\mathbb{R}^3 \times S^1 \quad \text{mit} \quad r_4 = L_0 = \xipar \cdot \lP = 2{,}155 \times 10^{-39}\,\text{m}}
\end{equation}
\end{keyresult}


\begin{foundation}[Hinweis: Verbindung zum SI-System]
Die absoluten Zahlenwerte in diesem Dokument — insbesondere $L_0 = 2{,}155 \times 10^{-39}\,\text{m}$ — setzen eine Verbindung zum SI-Einheitensystem voraus. Diese ist nicht zirkulär, sondern epistemologisch notwendig: Aus einem dimensionslosen Parameter $\xipar$ allein lassen sich keine absoluten SI-Werte ableiten. Die Verbindung erfolgt über $m_e$, $\lP$ oder $\alpha$, die alle äquivalente Einstiegspunkte sind. Ausführlich behandelt in Dokument 180 (Abschnitt: Verbindung zum SI-System) sowie in Dokument 056 und 101.

Dokumente 180, 181 und 182 bilden eine Einheit und sollten nicht getrennt gelesen werden.
\end{foundation}




\end{document}
